\origpage{66}{ROBERTSON AND SEYMOUR}
\noindent there is a $(p+4)$-nest $(C_1, \dots, C_{p+4})$ and disjoint $X \to \bar{\Omega}_0$ paths $P_1, \dots, P_{2p+8}$ of $G'$, such that for $1 \le i \le 2p+8$ and $1 \le j \le p+4$, $P_i \cap C_j$ is a path. Since every vertex of $G'$ in $\bar{\Omega}$ has valency $\le 1$ it follows that for $1 \le j \le p+4$ no vertex of $C_j$ is in $\bar{\Omega}$ and that for $1 \le i \le 2p+8$ no vertex of $P_i$ is in $\bar{\Omega}$ except its initial vertex. Let $x_i$ be the initial vertex of $P_i$ and let $y_i$ be the first vertex of $P_i$ in $V(C_{p+4})$ ($1 \le i \le 2p+8$). Put $Y = \{y_1, \dots, y_{2p+8}\}$. Let $G_1$ be the subgraph of $G$ drawn in the closed disc $\Delta_1$ bounded by the drawing of $C_{p+4}$ (thus, $C_{p+4} \subseteq G_1$). Let $\Omega_1$ be the cyclic permutation in $V(G_1)$ with $\bar{\Omega}_1 = V(C_{p+4})$ corresponding to the clockwise orientation of $bd(\Delta_1)$. Then $(G \cup G_0, \Omega)$ is the composition of $(G_1 \cup G_0, \Omega_1)$ with some rural neighborhood, and so by (9.1), $(G_1 \cup G_0, \Omega_1)$ has a crooked transaction of cardinality $\ge p$. But $Y \subseteq \bar{\Omega}_1$ is weakly-linked in the rural neighborhood $(G_1, \Omega_1, \Omega_0)$, because our choices of $C_1, \dots, C_{p+4}$ and $P_1, \dots, P_{2p+8}$ ensure that for $1 \le i \le 2p+8$ the subpath of $P_i$ from $Y$ to $\bar{\Omega}_0$ is a path of $G_1$. This neighborhood has breadth $\ge p+4$ and so by (10.3), there is a crooked transaction $\mathscr{P}$ in $(G_1 \cup G_0, \Omega_1)$ of cardinality $\ge p$ with $V(\mathscr{P}) \subseteq Y$. For each $P \in \mathscr{P}$, let its ends be $y_i, y_{i'} \in Y$, where $1 \le i, i' \le 2p+8$, and let $P^+$ denote the union of $P$ with the subpath of $P_i$ between $x_i$ and $y_i$ and the subpath of $P_{i'}$ between $x_{i'}$ and $y_{i'}$. Then $P^+$ is a path of $G\fix{ \cup G_0}$\fixnote{fix:10.4-path}{Here $P$ is a path of $G_1 \cup G_0$ and may use edges of $G_0$, so $P^+$ is a path of $G \cup G_0$, as a member of a transaction in $(G \cup G_0, \Omega)$ should be.} with both ends in $X$ and no internal vertex in $\bar{\Omega}$, and the set $\{P^+ : P \in \mathscr{P}\}$ is the desired crooked transaction. \fix{We check this. The subpath of $P_i$ between $x_i$ and $y_i$ meets $\Delta_1$ only in $y_i$, since
$y_i$ is the first vertex of $P_i$ in $V(C_{p+4})$ and $x_i \in bd(\Delta)$ is not in $\Delta_1$. So
this subpath meets $G_1 \cup G_0$ only in $y_i$, because $G_1$ is drawn in $\Delta_1$ and
$V(G \cap G_0) = \bar{\Omega}_0$ is drawn in $\Delta_0 \subseteq \Delta_1$. Since also the members of
$\mathscr{P}$ are disjoint and $P_1, \dots, P_{2p+8}$ are disjoint, the paths $P^+$
$(P \in \mathscr{P})$ are mutually disjoint. Let $H_2$ be the subgraph of $G$ drawn in the closure of
$\Delta - \Delta_1$. Then $(H_2, \Omega, \Omega_1)$ is a rural neighborhood, and the subpaths of
$P_1, \dots, P_{2p+8}$ between $x_i$ and $y_i$ are disjoint paths of $H_2$ from $\bar{\Omega}$ to
$\bar{\Omega}_1$. Hence for distinct $i_1, i_2, i_3, i_4$, $\{x_{i_1}, x_{i_3}\}$ crosses
$\{x_{i_2}, x_{i_4}\}$ in $\Omega$ if and only if $\{y_{i_1}, y_{i_3}\}$ crosses
$\{y_{i_2}, y_{i_4}\}$ in $\Omega_1$. A cyclic order of a finite set is determined, up to reversal, by
which of its pairs cross. So the ends of the paths $P^+$ lie around $\Omega$ in the same cyclic order,
up to reversal, as the corresponding ends of the members of $\mathscr{P}$ around $\Omega_1$. Whether a
set of disjoint paths is a transaction, and which of its members are peripheral, depends only on
this order. Hence $\{P^+ : P \in \mathscr{P}\}$ is a crooked transaction in $(G \cup G_0, \Omega)$,
of cardinality $|\mathscr{P}| \ge p$.\fixnote{fix:10.4-proof-added}{The paper says that the new
paths form a crooked transaction and gives no reason. The reason above is the one used in the
formalization.}} This completes the proof. $\square$

In later papers we shall also require the following, which we include here because of its similarity to (10.4).

\stmt
(10.5) \textit{Let $p \ge 0$ and let $(G, \Omega, \Omega_0)$ be a rural neighborhood of breadth $\ge \frac{1}{2} p + 1$. Let $X \subseteq \bar{\Omega}$ be linked, with $|X| \ge p$. Let $(G \cup G_0, \Omega)$ be the composition of some society $(G_0, \Omega_0)$ with $(G, \Omega, \Omega_0)$, and let $Y \subseteq V(G_0)$. Then either there are $p$ mutually disjoint paths of $G \cup G_0$ from $X$ to $Y$, each with only its initial vertex in $\bar{\Omega}$, or there exists $W \subseteq V(G \cup G_0)$ with $|W| < p$ and $W \cap \bar{\Omega} = \emptyset$ such that every path of $G \cup G_0$ from $\bar{\Omega}$ to $Y$ has a vertex in $W$.}

\textit{Proof.} Suppose that there is no $W$ as in the theorem. Let $G' = (G \cup G_0) \setminus (\bar{\Omega} - X)$. It suffices to show that there are $p$ mutually disjoint $X \to Y$ paths in $G'$; and therefore, by Menger's theorem, it suffices to show that if $Z \subseteq V(G')$ and $|Z| < p$ then there is a path of $G'$ from $X$ to $Y$ avoiding $Z$. Thus, let $Z \subseteq V(G')$ with $|Z| < p$. Let $s = \lceil \frac{1}{2} p \rceil$; as in (10.4), there is an $s$-nest $(C_1, \dots, C_s)$ in some presentation of $(G, \Omega, \Omega_0)$, such that $C_1, \dots, C_s \subseteq G'$ and $V(C_s) \cap \bar{\Omega} = \emptyset$. Since $X$ is linked and $|X| \ge p$, there is a path $P_1$ of $G'$ from $X$ to $\bar{\Omega}_0$ which avoids $Z$. Since $C_1, \dots, C_s$ are mutually disjoint, there exists $j$ with $1 \le j \le s$ such that $|Z \cap V(C_j)| \le 1$. Since there is no $W$ as in the theorem, there are $p$ paths of $G \cup G_0$ from $\bar{\Omega}$ to $Y$, mutually vertex-disjoint except possibly for their initial vertices, by a form of Menger's theorem. Since these all meet $V(C_j)$, there is a path $P_2$ of $G'$
